\section{Mathematical Formulation}
\label{sec:implementation}

\subsection{Weisfeiler-Lehman Graph Kernel}
\label{subsec:wl_kernel}

The Weisfeiler-Lehman subtree kernel~\cite{ref_shervashidze2011} measures structural similarity by iteratively re-labeling each node via compressed hashing to encode the node label and neighbors up to hop distance $h$. This forms a feature map of local subgraph topology independent of gradient descent training. Phase~1 demonstrated $h=2$ to be the optimal depth for method-level CS-1 program graphs; an ablation analysis of depth (Table~\ref{tab:wl_ablation}) using the 400-program cross-language dataset proves the choice correct: $h=1$ undersamples method-level encapsulation (85.1\%), and $h=3$ oversamples due to hash fragmentation from greater subgraph dissimilarity at language frontend boundaries (87.0\%), representing a -0.25~pp regression from the optimal $h=2$ score of 87.25\%.

Once the graph has been hashed via WL kernel, each program can be mapped into a frequency vector $\phi_{WL}(G)$. Structural similarity between a query graph $q$ and reference graph $r$ is given by:
\begin{equation}
    s_{WL}(q,r) = \frac{\langle \phi_{WL}(G_q),\, \phi_{WL}(G_r) \rangle}{\|\phi_{WL}(G_q)\| \cdot \|\phi_{WL}(G_r)\|}
\end{equation}

\textbf{Computational complexity:} The WL subtree kernel runs in $O(h \cdot m)$ time per graph, where $h$ is the number of refinement steps ($h{=}2$) and $m$ is the number of edges in the CPG. For introductory CS-1 programs, $m$ typically ranges between 100 and 400 edges. Pairwise structural comparison between $N$ student submissions and $R$ reference solutions runs in $O(N \cdot R \cdot |V|)$ time, where $|V|$ is the size of the combined WL feature vocabulary. Tversky similarity for the CES taxonomy runs in $O(|P|)$ time, where $|P|{=}21$ is the fixed number of patterns in the taxonomy. Late fusion combination runs in $O(1)$ time per program pair. The overall execution cost is dominated by the Joern CPG extraction process, which runs in constant time relative to the small size of method-level CS-1 assignments, yielding an empirical average of ${\sim}1.2$~s per submission (Section~\ref{sec:results}).

\subsection{Extended CES Similarity with Asymmetric Tversky Index}
\label{subsec:tversky}

The Tversky similarity~\cite{ref_tversky1977} offers a principled generalization of Jaccard similarity in the asymmetrical case. The set representations of $A$ (student submission) and $B$ (reference solution) use the call-depth importance function defined below to compute weights $w_{p}$ for each unique pattern $p$, and collapses multiple instances of a particular pattern type in a single program to their highest weight when comparing submissions. The Tversky similarity is:
\begin{equation}
    T_{\alpha,\beta}(A,B) = \frac{I(A,B)}{I(A,B) + \alpha\,F(A,B) + \beta\,M(A,B)}
\end{equation}
where $I(A,B) = \sum_{p \in A \cap B} w_p$ is the weighted intersection of shared pattern types, $F(A,B) = \sum_{p \in A \setminus B} w_p$ represents student-only patterns absent from the reference, and $M(A,B) = \sum_{p \in B \setminus A} w_p$ represents reference patterns absent from the submission. Setting $\alpha{=}0.1$ minimally punishes student-only patterns while $\beta{=}0.9$ harshly punishes lack of reference patterns---this reflects the notion in education that the correct solution should include all key computational patterns found in the reference, notwithstanding additional syntactic fluff. These values offset the verbosity differential of the Java language relative to C/C++, resulting in spurious CES patterns in Java submissions.

These values of $\alpha{=}0.1$ and $\beta{=}0.9$ were fixed before evaluating the approach based on the educational insight above. A range study over $\alpha \in \{0.05, 0.10, 0.15, 0.20\}$ demonstrates consistency at $\pm0.1\%$ variance for the monolingual group and $\pm0.35\%$ variation for the cross-language corpus (Table~\ref{tab:tversky_sensitivity}).

Weight assignments for the CES pattern are done based on a call-depth-based importance: the outermost containing function is assigned a weight of 1.0, its first-degree callees get a weight of 0.8, second-degree callees 0.7, and third-degree or further callees 0.6. The ablation study on the monolingual dataset in C++ shows that assigning equal call-depth-based weights leads to deterioration of top-1 recall performance from 90.0\% to 88.5\%, thus allowing us to attribute the role of the depth-decay mechanism in reducing the noise of low-level utility functions. This ablation was conducted on the C++ monolingual cohort only; a corresponding ablation on the cross-language cohort is deferred to future work due to dataset size constraints and is documented as a validity boundary in Section~\ref{subsec:threats}.

\begin{table}[ht]
\caption{WL Depth Ablation on Cross-Language Dataset ($N=400$, 95\% Wilson CI).}
\label{tab:wl_ablation}
\centering
\begin{tabular}{|c|c|c|}
\hline
\textbf{WL Depth} & \textbf{Top-1 Accuracy} & \textbf{95\% Wilson CI} \\
\hline
$h=1$ & 85.10\% & [81.26\%, 88.31\%] \\
$h=2$ & \textbf{87.25\%} & [83.60\%, 90.20\%] \\
$h=3$ & 87.00\% & [83.33\%, 90.00\%] \\
\hline
\end{tabular}
\end{table}

\begin{table}[ht]
\caption{Tversky Parameter Sensitivity Analysis. Accuracy results presented for monolingual C++ ($N=400$) and cross-language ($N=400$) dataset samples. The chosen combination ($\alpha{=}0.1$, $\beta{=}0.9$) yields maximum accuracy for both populations of students.}
\label{tab:tversky_sensitivity}
\centering
\begin{tabular}{|c|c|c|c|}
\hline
$\boldsymbol{\alpha}$ & $\boldsymbol{\beta}$ & \textbf{C++ Acc.} & \textbf{Cross-Lang. Acc.} \\
\hline
0.05 & 0.95 & 89.75\% & 86.90\% \\
0.10 & 0.90 & \textbf{90.00\%} & \textbf{87.25\%} \\
0.15 & 0.85 & 89.75\% & 87.25\% \\
0.20 & 0.80 & 89.50\% & 86.75\% \\
\hline
\end{tabular}
\end{table}

\subsection{Late-Fusion Score and Hard-Filtering Constraint}
\label{subsec:late_fusion}

The three view scores are combined via linear weighted late fusion. For a student query submission $s_t$ and reference candidate $r$, the final similarity score is:
\begin{equation}
\begin{split}
    s_{\text{final}}(s_t, r) = &\; w_{\text{BL}} \cdot s_{\text{BL}}(s_t, r) \\
                               +&\; w_{\text{WL}} \cdot s_{\text{WL}}(s_t, r) \\
                               +&\; w_{\text{CES}} \cdot s_{\text{CES}}(s_t, r)
\end{split}
\end{equation}
where $(w_{\text{BL}}, w_{\text{WL}}, w_{\text{CES}}) = (0.35, 0.40, 0.25)$ are the Phase~1 weights applied without modification. Top-1 retrieval selects the reference $r^* = \arg\max_{r \in \mathcal{R}} s_{\text{final}}(s_t, r)$.

\textbf{Strict Hard-Filtering Constraint:} To prevent trivial same-language retrieval from inflating cross-language accuracy, the reference set $\mathcal{R}$ is pruned to exclude all programs written in the same language as the query:
\begin{equation}
    \mathcal{R}_{\text{filtered}}(s_t) = \{r \in \mathcal{R} \mid \text{lang}(r) \neq \text{lang}(s_t)\}
\end{equation}
Cross-language reference pairs (e.g., a C query scored against a C++ reference) remain eligible; only same-language mappings are blocked. This constraint ensures that all reported accuracy figures reflect genuine cross-language retrieval performance.

\subsection{Weight Transfer Stability Under Cross-Language Evaluation}
\label{subsec:weight_stability}

Using the Phase~1 optimal weights $(0.35, 0.40, 0.25)$---trained on 104 C-language programs---without re-training resulted in competitive top-1 retrieval performance under the current cross-language setting (87.25\% top-1 accuracy). Given that the vast majority of retrieval pairs derive from the common C/C++ front-end, this should be considered partial transfer stability and not language-independence. Under the conditions of the current evaluation, no statistically significant deterioration occurred (two-proportion z-test, $z \approx 1.44$, $p \approx 0.15$; see Section~\ref{subsec:weight_stability_discussion} for complete analysis); lack of statistical significance does not mean equivalence under other language settings. Table~\ref{tab:weight_stability} lists the weights used in all evaluations.
\begin{table}[ht]
\caption{Weight Configurations and Retrieval Accuracy Across Evaluation Contexts.}
\label{tab:weight_stability}
\centering
\begin{tabular}{|l|c|c|c|c|}
\hline
\textbf{Context} & $\boldsymbol{w_{\text{BL}}}$ & $\boldsymbol{w_{\text{WL}}}$ & $\boldsymbol{w_{\text{CES}}}$ & \textbf{Top-1 Acc.} \\
\hline
Phase 1 C (source) & 0.35 & 0.40 & 0.25 & 88.50\% \\
C++ monolingual & 0.35 & 0.40 & 0.25 & 90.00\% \\
Java monolingual & 0.35 & 0.40 & 0.25 & 85.75\% \\
Cross-language & 0.35 & 0.40 & 0.25 & 87.25\% \\
C++ in-sample opt.$^*$ & 0.00 & 0.05 & 0.95 & 90.50\% \\
\hline
\end{tabular}

\vspace{4pt}
{\raggedright \footnotesize $^*$ C++ in-sample optimum; not used in cross-language evaluation.\par}
\vspace{-1em} % Adjust this value (e.g., -1.5em or -10pt) to pull the text up more or less
\end{table}